% 2.总体分析.tex —— 二、总体分析
\section{总体分析}
四个问题采用相同的径向几何、初始场和表面换热传质边界，并依次扩展物性与输出要求。一维近似描述长圆柱中段，不直接验证有限长度药材的端部状态；问题三、四的“各处达标”据此解释为所建径向模型内的全场达标。

\textbf{问题一：先明确物理依赖，再比较数值格式。}附录 2 中热物性取常数，热方程不含 $C$；水分扩散系数仅为 $D(C)$，不含 $T$。因此两个方程\textbf{相互解耦}，程序中的先热后质只是计算顺序。空间上分别处理中心奇异项、内部通量与表面 Robin 条件；时间上先建立显式 Euler 基线，再用 CN 和 Thomas 追赶法求解。CN 对线性扩散离散系统具有无条件稳定性，但这不同时保证大步长下的精度、单调性或非线性迭代收敛。

\textbf{问题二：通过变物性描述全程耦合。}$\rho(C),c_p(C),k(C)$ 使含水率影响温度场，$D(C,T)$ 则使温度影响水分迁移。热方程采用局部等效热容近似，未建立包含物质迁移焓与汽化潜热的完整能量模型。数值上，在水分块中保留 $\partial D/\partial C$ 敏度项，并与热方程交替更新。这是含 Newton 型线性化的块交替求解，局部迭代表现不能推导为整个非线性系统的二次收敛保证。环境在前 $4$ h 按附件插值、之后按末值延拓，模拟至 $72$ h；药材提前接近环境温度不改变所设边界阶段。

\textbf{问题三：分开处理终止判据与离散误差。}在问题二的全程模型上，使用全节点最大含水率判断首次达标，以时间步加倍估计控制推进，并在越阈区间重新积分、二分定位。输出插值误差、时间步误差和空间网格误差分别检查；均匀与分级网格比较用于估计终止时间的网格依赖。准稳态形状漂移由同一数值轨道的后续数据拟合，只用于事后降阶一致性分析。

\textbf{问题四：以移动边界刻画收缩。}问题一至三假设半径固定，问题四引入附件 2 的收缩半径 $R(t)$ 作为移动边界，采用 Landau 变换把移动域 $[0,R(t)]$ 映射为固定域 $[0,1]$，由链式法则产生对流项 $+(E R'/R)\partial C/\partial E$（题面符号为笔误）。收缩高度前重、集中在前 $24$ h，故对流项只在早期活跃；后半程退化为定半径干燥。烘干时长同时做空间与时间收敛检验，并给出有/无收缩、不同对流项符号的对比以定位收缩效应与符号正误。

图表依次服务于上述论证：问题一展示径向剖面与格式差异，问题二展示物性变化和全程剖面，问题三展示时间步控制、空间加密及全场达标过程，问题四展示收缩规律、移动边界剖面与收缩对烘干时长的缩短。每幅图的\textbf{验证对象}在图注中注明，分别指向模型参数、数值格式、物理规律或最终结论；每幅图的结论限定于其数据和计算设置，避免用稳定性图代替精度检验，或用同轨道拟合代替独立验证。
